% Written by Claude (2026-10-14): initialization, first of three
% training-in-practice slides. Figure: shared/assets/images/lecture_4/init/.
\begin{frame}{Initializing the weights}
	\MLPTrainingFont
	\setlength{\abovedisplayskip}{4pt}\setlength{\belowdisplayskip}{4pt}
	\begin{columns}[T,onlytextwidth]
		\begin{column}{0.50\textwidth}
			\begin{tcolorbox}[nndsorange,boxsep=1mm,top=1mm,bottom=1mm]
				If all the weights of a layer start equal, all its hidden units stay identical during training. We break this \myalert{symmetry} with random weights.
			\end{tcolorbox}

			The scale also matters. With weights of variance $\sigma^2$, each ReLU layer multiplies the mean square of the activations by $\sigma^2 D_{\ell-1}/2$.

			\begin{tcolorbox}[nndsblue,boxsep=1mm,top=1mm,bottom=1mm]
				\textbf{\color{darkcerulean}He initialization}
				\[
					W^{(\ell)}_{ij}\sim\mathcal{N}\!\left(0,\frac{2}{D_{\ell-1}}\right),
					\qquad \vc{b}^{(\ell)}=\vc{0}.
				\]
			\end{tcolorbox}
		\end{column}
		\begin{column}{0.46\textwidth}
			\centering
			\textbf{\color{darkcerulean}Root mean square of $\vc{h}^{(\ell)}$}\par
			{\scriptsize\color{drawgray}ReLU MLP, width $D_\ell=512$}\par
			\medskip
			% Generated by plot_activation_scale.py; do not edit by hand.
\begin{tikzpicture}
  \draw[drawgray!35] (0,0) -- (4.00,0);
  \draw[->,drawgray] (0,-1.620) -- (0,1.755);
  \draw[->,drawgray] (0,-1.620) -- (4.30,-1.620);
  \node[font=\scriptsize] at (2.00,-2.170) {layer $\ell$};
  \draw[drawgray] (0.00,-1.620) -- ++(0,-0.06) node[below,font=\scriptsize] {$0$};
  \draw[drawgray] (1.00,-1.620) -- ++(0,-0.06) node[below,font=\scriptsize] {$5$};
  \draw[drawgray] (2.00,-1.620) -- ++(0,-0.06) node[below,font=\scriptsize] {$10$};
  \draw[drawgray] (3.00,-1.620) -- ++(0,-0.06) node[below,font=\scriptsize] {$15$};
  \draw[drawgray] (4.00,-1.620) -- ++(0,-0.06) node[below,font=\scriptsize] {$20$};
  \draw[drawgray] (0,-1.350) -- ++(-0.06,0) node[left,font=\scriptsize] {$2^{-10}$};
  \draw[drawgray] (0,-0.675) -- ++(-0.06,0) node[left,font=\scriptsize] {$2^{-5}$};
  \draw[drawgray] (0,0.000) -- ++(-0.06,0) node[left,font=\scriptsize] {$1$};
  \draw[drawgray] (0,0.675) -- ++(-0.06,0) node[left,font=\scriptsize] {$2^{5}$};
  \draw[drawgray] (0,1.350) -- ++(-0.06,0) node[left,font=\scriptsize] {$2^{10}$};
  \draw[drawred,very thick,densely dashed] plot coordinates {(0.000,0.000) (0.200,-0.067) (0.400,-0.126) (0.600,-0.188) (0.800,-0.259) (1.000,-0.336) (1.200,-0.410) (1.400,-0.476) (1.600,-0.551) (1.800,-0.627) (2.000,-0.707) (2.200,-0.772) (2.400,-0.855) (2.600,-0.926) (2.800,-1.009) (3.000,-1.065) (3.200,-1.122) (3.400,-1.192) (3.600,-1.263) (3.800,-1.329) (4.000,-1.395)};
  \node[right,text=drawred,font=\scriptsize] at (4.00,-1.395) {$\sigma^2=1/D$};
  \draw[darkcerulean,very thick,solid] plot coordinates {(0.000,0.000) (0.200,-0.001) (0.400,-0.008) (0.600,-0.013) (0.800,-0.010) (1.000,-0.012) (1.200,0.003) (1.400,-0.005) (1.600,-0.011) (1.800,-0.006) (2.000,-0.012) (2.200,-0.024) (2.400,-0.026) (2.600,-0.027) (2.800,-0.029) (3.000,-0.035) (3.200,-0.050) (3.400,-0.065) (3.600,-0.058) (3.800,-0.067) (4.000,-0.071)};
  \node[right,text=darkcerulean,font=\scriptsize] at (4.00,-0.071) {$\sigma^2=2/D$};
  \draw[draworange,very thick,densely dotted] plot coordinates {(0.000,0.000) (0.200,0.068) (0.400,0.131) (0.600,0.191) (0.800,0.277) (1.000,0.347) (1.200,0.414) (1.400,0.472) (1.600,0.540) (1.800,0.606) (2.000,0.667) (2.200,0.743) (2.400,0.804) (2.600,0.877) (2.800,0.963) (3.000,1.034) (3.200,1.095) (3.400,1.176) (3.600,1.244) (3.800,1.310) (4.000,1.384)};
  \node[right,text=draworange,font=\scriptsize] at (4.00,1.384) {$\sigma^2=4/D$};
\end{tikzpicture}

		\end{column}
	\end{columns}
	\vfill
	{\scriptsize\color{drawgray}The same reasoning applies to gradients flowing backward, as we will see in the autodiff lecture. For other activations, see \href{https://proceedings.mlr.press/v9/glorot10a.html}{Glorot and Bengio (2010)}; \href{https://arxiv.org/abs/1502.01852}{He et al. (2015)}.\par}
\end{frame}
